% nyenje_report_clean.tex
\documentclass[12pt]{article}
\usepackage[utf8]{inputenc}
\usepackage{geometry}
\geometry{a4paper, margin=1in}
\usepackage{amsmath}
\usepackage{amsfonts}
\usepackage{booktabs}
\usepackage{graphicx}
\usepackage{hyperref}
\usepackage{xcolor}

% Define colors
\definecolor{primary}{RGB}{30,60,114}
\definecolor{secondary}{RGB}{42,82,152}

\begin{document}

% Title Page
\begin{titlepage}
    \centering
    \vspace*{2cm}
    
    {\huge\bfseries\color{primary} Nyenje Bay Wind/Waves Analysis Tool}\\[0.5cm]
    {\Large Technical Framework Report}\\[2cm]
    
    {\large Prepared by: \textbf{Key Informatics}}\\[0.3cm]
    {\large Location: \textbf{Nyenje Bay, Lake Kariba}}\\[0.3cm]
    {\large Coordinates: \textbf{16.53S, 28.83E}}\\[0.3cm]
    {\large Date: \textbf{April 2026}}\\[2cm]
    
    \vfill
    
    {\large \textit{Engineering Design Basis for Floating Photovoltaic (FPV) Structures}}
\end{titlepage}

% Table of Contents
\tableofcontents
\newpage

% Executive Summary
\section*{Executive Summary}
\addcontentsline{toc}{section}{Executive Summary}

This report documents the technical framework of the Nyenje Bay Wind/Waves Analysis Tool. The tool integrates multiple layers:
\begin{itemize}
    \item \textbf{User Interface Layer:} Streamlit-based web dashboard
    \item \textbf{Scientific Computing Layer:} NumPy, SciPy, Pandas, Xarray
    \item \textbf{Data Layer:} ERA5 reanalysis, CCMP satellite, ground measurements
    \item \textbf{Wave Modeling Layer:} SWAN model integration
\end{itemize}

\vspace{0.5cm}

\textbf{Key Scientific Calculations:}
\begin{itemize}
    \item Wind averaging periods: 3-minute, 10-minute, 3-second gust, 10-second gust
    \item Gust factors: 1.15 (3-min), 1.45 (3-sec), 1.35 (10-sec)
    \item JONSWAP wave height formula
    \item Gumbel extreme value analysis for return periods
\end{itemize}

\newpage

% Section 1
\section{Tool Architecture}

\subsection{Four-Layer Framework}

The tool is built on a four-layer architecture:

\begin{table}[h]
\centering
\begin{tabular}{|p{3cm}|p{4cm}|p{6cm}|}
\hline
\textbf{Layer} & \textbf{Technology} & \textbf{Functions} \\
\hline
User Interface & Streamlit & Web dashboard, interactive plots \\
Scientific Computing & NumPy, SciPy & Wind averaging, wave calculations \\
Data Processing & Xarray, Pandas & ERA5 reading, grid interpolation \\
External Models & SWAN & Wave distribution \\
\hline
\end{tabular}
\caption{Architecture Layers}
\end{table}

\newpage

% Section 2
\section{Wind Averaging Calculations}

\subsection{Overview}

The tool calculates five different wind averaging periods:

\begin{table}[h]
\centering
\begin{tabular}{|l|l|l|}
\hline
\textbf{Period} & \textbf{Gust Factor} & \textbf{Application} \\
\hline
Hourly Mean & 1.00 & Reference wind speed \\
3-Minute Average & 1.15 & Dynamic response \\
10-Minute Average & 1.00 & Marine forecast \\
3-Second Gust & 1.45 & Ultimate Limit State (ULS) \\
10-Second Gust & 1.35 & Serviceability Limit State (SLS) \\
\hline
\end{tabular}
\caption{Wind Averaging Periods}
\end{table}

\subsection{Exponential Weighted Moving Average}

For hourly data, the tool uses Exponential Weighted Moving Average:

\begin{equation}
\text{EWMA}_t = \alpha \cdot x_t + (1 - \alpha) \cdot \text{EWMA}_{t-1}
\end{equation}

Where:
\begin{itemize}
    \item $\alpha = \frac{2}{\text{span} + 1}$ is the smoothing factor
    \item For 3-minute average: span = 3, $\alpha = 0.5$
    \item For 10-minute average: span = 10, $\alpha \approx 0.18$
\end{itemize}

\subsection{Gust Factor Formula}

\begin{equation}
\text{Gust Factor} = \frac{\text{Peak Gust Speed}}{\text{Mean Wind Speed}}
\end{equation}

\begin{table}[h]
\centering
\begin{tabular}{|l|l|}
\hline
\textbf{Period} & \textbf{Gust Factor} \\
\hline
3-second gust & 1.45 \\
10-second gust & 1.35 \\
3-minute average & 1.15 \\
\hline
\end{tabular}
\caption{Gust Factor Derivation}
\end{table}

\newpage

% Section 3
\section{Wave Height Calculation}

\subsection{JONSWAP Formula}

The Joint North Sea Wave Project (JONSWAP) formula is used for fetch-limited waves.

\subsubsection{Dimensionless Fetch}

\begin{equation}
\tilde{x} = \frac{g \cdot F}{U^2}
\end{equation}

Where:
\begin{itemize}
    \item $g = 9.81 \text{ m/s}^2$ (gravity)
    \item $F$ = fetch length (meters)
    \item $U$ = wind speed (m/s)
\end{itemize}

\subsubsection{Significant Wave Height}

\begin{equation}
H_s = 0.0016 \cdot \frac{U^2}{g} \cdot \sqrt{\tilde{x}}
\end{equation}

\subsubsection{Peak Wave Period}

\begin{equation}
T_p = 0.286 \cdot \frac{U}{g} \cdot (\tilde{x})^{0.33}
\end{equation}

\subsection{Simplified Formula}

\begin{equation}
\boxed{H_s = 0.0016 \cdot U \cdot \sqrt{g \cdot F}}
\end{equation}

\subsection{Example Calculations}

\begin{table}[h]
\centering
\begin{tabular}{|l|l|l|}
\hline
\textbf{Wind Speed (m/s)} & \textbf{Wave Height Hs (m)} & \textbf{Period Tp (s)} \\
\hline
5 & 0.07 & 1.2 \\
8 & 0.18 & 2.0 \\
10 & 0.29 & 2.3 \\
12 & 0.41 & 2.7 \\
15 & 0.64 & 3.2 \\
18 & 0.92 & 3.7 \\
20 & 1.14 & 4.0 \\
\hline
\end{tabular}
\caption{Wave Height for Different Wind Speeds (Fetch = 4 km)}
\end{table}

\newpage

% Section 4
\section{Extreme Value Analysis}

\subsection{Gumbel Distribution}

The cumulative distribution function (CDF) is:

\begin{equation}
F(x) = \exp\left(-\exp\left(-\frac{x - \mu}{\beta}\right)\right)
\end{equation}

Where:
\begin{itemize}
    \item $\mu$ = location parameter
    \item $\beta$ = scale parameter
\end{itemize}

\subsection{Return Period Calculation}

For return period T (years):

\begin{equation}
x(T) = \mu - \beta \cdot \ln\left(-\ln\left(1 - \frac{1}{T}\right)\right)
\end{equation}

\subsection{Results for Nyenje Bay}

\begin{table}[h]
\centering
\begin{tabular}{|l|l|l|}
\hline
\textbf{Return Period (years)} & \textbf{Wind Speed (m/s)} & \textbf{Wave Height (m)} \\
\hline
2 & 11.2 & 0.85 \\
5 & 12.8 & 1.02 \\
10 & 14.0 & 1.15 \\
20 & 15.2 & 1.28 \\
25 & 15.6 & 1.32 \\
50 & 16.8 & 1.48 \\
100 & 18.0 & 1.62 \\
\hline
\end{tabular}
\caption{Return Period Wave Heights}
\end{table}

\newpage

% Section 5
\section{Validation Results}

\subsection{Kariba Airport Comparison}

\begin{table}[h]
\centering
\begin{tabular}{|l|l|l|l|}
\hline
\textbf{Metric} & \textbf{Raw ERA5} & \textbf{Debiased ERA5} & \textbf{Improvement} \\
\hline
Correlation (r) & 0.39 & 0.502 & +29\% \\
Bias (m/s) & +2.04 & +0.01 & 99.5\% \\
RMSE (m/s) & 2.45 & 0.52 & 79\% \\
\hline
\end{tabular}
\caption{Validation Statistics (2001-2005)}
\end{table}

\subsection{Interpretation of r = 0.502}

\begin{equation}
r = 0.502 \quad \Rightarrow \quad R^2 = 0.252
\end{equation}

This means:
\begin{itemize}
    \item 25.2\% of wind variance is shared between locations
    \item 74.8\% is due to local effects (expected for 3 km separation)
    \item The correlation exceeds the expected value of 0.40
\end{itemize}

\newpage

% Section 6
\section{Technology Stack}

\begin{table}[h]
\centering
\begin{tabular}{|l|l|l|}
\hline
\textbf{Component} & \textbf{Technology} & \textbf{Purpose} \\
\hline
UI Framework & Streamlit & Web dashboard \\
Scientific & NumPy & Array operations \\
Data Analysis & Pandas & DataFrames \\
NetCDF & Xarray & ERA5 reading \\
Statistics & SciPy & Gumbel distribution \\
Visualization & Plotly & Interactive plots \\
Wave Model & SWAN & Wave distribution \\
\hline
\end{tabular}
\caption{Technology Stack}
\end{table}

\newpage

% Section 7
\section{Conclusions}

\subsection{Key Technical Achievements}

\begin{enumerate}
    \item \textbf{Multi-Layer Architecture}: Separation of UI, computation, and data
    \item \textbf{Wind Averaging}: EWMA method for 3-min and 10-min averages
    \item \textbf{Gust Factors}: Engineering-standard factors (1.45, 1.35)
    \item \textbf{Wave Calculations}: JONSWAP fetch-limited formula
    \item \textbf{Extreme Analysis}: Gumbel distribution for return periods
    \item \textbf{Validation}: r = 0.502 correlation with ground measurements
\end{enumerate}

\subsection{Framework Independence}

The scientific core is completely independent of the user interface:
\begin{itemize}
    \item Same calculations work in command line, web app, or API
    \item Streamlit is only for visualization, not computation
    \item Core functions can be reused in other projects
\end{itemize}

\end{document}